% Generated draft: requires mathematical and source review.
\documentclass[11pt]{article}
\usepackage[margin=25mm]{geometry}
\usepackage{fontspec}
\setmainfont{Latin Modern Roman}
\usepackage{amsmath,amssymb,mathtools,unicode-math}
\setmathfont{Latin Modern Math}
\usepackage{enumitem,needspace,xurl,hyperref}
\setlength{\parindent}{0pt}
\setlength{\parskip}{5pt}
\setlength{\emergencystretch}{3em}
\newcommand{\N}{\mathbb{N}}
\newcommand{\Z}{\mathbb{Z}}
\newcommand{\Q}{\mathbb{Q}}
\newcommand{\R}{\mathbb{R}}
\newcommand{\C}{\mathbb{C}}
\newcommand{\F}{\mathbb{F}}
\newcommand{\A}{\mathbb{A}}
\newcommand{\PP}{\mathbb{P}}
\newcommand{\E}{\mathbb{E}}
\newcommand{\eps}{\varepsilon}
\newcommand{\dd}{\,\mathrm{d}}
\begin{document}
\section*{UnsolvedMath ID: 2000011}
\textbf{Source problem code:} catalogue entry
\subsection{Definitions and mathematical statement}
Let \(\mathcal E_{(2.31)}\) denote the fracture-energy functional referenced as equation (2.31) in Ball's source. Its formula is not included in the supplied record. The problem asks to clarify the mathematical and physical status of models based on that specific functional, in relation to classical fracture theory and nonlinear elastostatics.

A variational fracture model typically specifies deformation variables, admissible cracks or displacement discontinuities, an energy, boundary/loading conditions, and a rule selecting equilibria or evolutions. These data distinguish static minimization, local stability, and irreversible crack growth; they cannot be inferred from an equation number alone. Consequently no particular bulk-plus-surface formula, small-strain approximation, finite-strain constitutive law, admissible function space, or evolution law is assigned to \(\mathcal E_{(2.31)}\) here. Without the omitted formula and its surrounding hypotheses, the requested comparison cannot be turned into a unique theorem or equivalence statement.

The supplied PARTIAL-PROGRESS summary describes substantial clarification since the 2000s for variational Griffith/Mumford--Shah fracture models, naming Francfort--Marigo and special functions of bounded variation (SBV), and reporting rigorous existence and phase-field \(\Gamma\)-convergence results. It says extensions to finite-strain nonlinear elastostatics have only partial results and fully general finite-elasticity fracture remains partly open. Those broad statements are retained as supplied survey assessments, not as universally quantified existence or convergence theorems for the absent functional.
\subsection{Short English statement}
Explain how the source's fracture-energy models relate to classical fracture and finite-strain elastic equilibrium. The defining equation (2.31) and its hypotheses are absent, so the exact model and comparison target require review.
\subsection{Sources}
\begin{itemize}
\item Local UnsolvedMath registry, permanent ID \(2000011\), exact source code \texttt{catalogue entry}, ``Some Open Problems in Elasticity --- Variational fracture models.''
\item Ball, elasticity problem list (2002), Problem 11, referencing equation (2.31); \url{https://people.maths.ox.ac.uk/ball/Articles%20in%20Conference%20Proceedings%20and%20Books/JMB%202002%20re%20Marsden%2060th.pdf}. PDF-text extraction; recorded access 2026-07-29. The equation itself is not supplied in the record.
\item Supplied research summary names Griffith/Mumford--Shah models, Francfort--Marigo, SBV existence theory, and phase-field \(\Gamma\)-convergence, without complete authoritative citations or precise result hypotheses.
\end{itemize}
\end{document}
